% =====================================================================
%  Tutorial: construcción de una app web con vibe coding (caso Byte)
%  Formato APA 7 (clase apa7, modo estudiante) + biblatex-apa.
%
%  Compilar desde docs/tutorial/:
%      latexmk -xelatex main.tex
%  (o: xelatex main && biber main && xelatex main && xelatex main)
%
%  Las figuras salen de docs/capturas/ con:
%      uv run python docs/tutorial/preparar_figuras.py
% =====================================================================
\documentclass[stu,12pt,floatsintext]{apa7}

% ---------- Datos de la portada (completar) ----------
\newcommand{\autor}{[Nombre completo del estudiante]}
\newcommand{\programa}{[Programa académico], [Nombre de la universidad]}
\newcommand{\curso}{Hacking ético}
\newcommand{\docente}{[Nombre del docente]}
\newcommand{\fecha}{13 de noviembre de 2026}

% ---------- Idioma y fuentes ----------
\usepackage[spanish,es-tabla,es-noshorthands,es-nolists]{babel}
\usepackage{fontspec}
\setmainfont{Times New Roman}
\setmonofont{Consolas}[Scale=0.88]
\usepackage{csquotes}
\usepackage[style=apa,backend=biber]{biblatex}
\addbibresource{referencias.bib}
\DefineBibliographyStrings{spanish}{nodate={s.\,f\adddot}}

% ---------- Paquetes ----------
\usepackage{graphicx}
\usepackage{booktabs}
\usepackage{float}
\usepackage{tabularx}
\usepackage{array}
\usepackage{enumitem}
\usepackage{xcolor}
\usepackage{fvextra}
\usepackage[most]{tcolorbox}
\usepackage{tikz}
\usetikzlibrary{arrows.meta,positioning}
\usepackage{pgfplots}
\pgfplotsset{compat=1.18}
\usepackage{ragged2e}
\usepackage{xurl}
\hypersetup{hidelinks}

\graphicspath{{figuras/}}

% APA 7: texto alineado a la izquierda (sin justificar) con sangría de 0,5 in.
\setlength{\RaggedRightParindent}{0.5in}
% APA 7: sin división de palabras al final de línea.
\setlength{\RaggedRightRightskip}{0pt plus 1fil}
\hyphenpenalty=10000
\exhyphenpenalty=10000
\AtBeginDocument{\def\keywordname{Palabras clave}\def\notelabel{Nota}}

% Código y salidas de terminal: interlineado sencillo y cortes de línea.
\fvset{breaklines=true,breakanywhere=true,fontsize=\small,frame=single,
  rulecolor=\color{gray!60},framesep=6pt,tabsize=4,
  formatcom=\linespread{1}\selectfont}

% Fragmentos citados del prompt maestro.
\definecolor{promptborde}{HTML}{3B4B8C}
\definecolor{promptfondo}{HTML}{F2F4FB}
\newtcolorbox[auto counter]{fragmento}[2][]{%
  breakable, enhanced, colback=promptfondo, colframe=promptborde,
  boxrule=0.6pt, arc=2pt, left=6pt, right=6pt, top=4pt, bottom=4pt,
  before skip balanced=12pt, after skip balanced=12pt,
  fonttitle=\bfseries\small, coltitle=white,
  title={Fragmento \thetcbcounter\ del prompt maestro. #2},
  fontupper=\small, before upper={\linespread{1}\selectfont\setlength{\parskip}{0.35em}\raggedright},
  #1}

% Cuadro de "lección" del tutorial.
\newtcolorbox{leccion}{%
  breakable, enhanced, colback=gray!5, colframe=gray!45, boxrule=0.5pt, arc=2pt,
  fonttitle=\bfseries\small, coltitle=black, colbacktitle=gray!18,
  title={Lección para replicar}, fontupper=\small,
  before upper={\linespread{1}\selectfont\raggedright}}

% Figura en formato APA: número en negrita, título en cursiva arriba, nota abajo.
% \figura{archivo}{ancho}{título}{nota}{etiqueta}
\newcommand{\figura}[5]{%
  \begin{figure}[H]
    \caption{#3}\label{#5}
    \centering
    \fbox{\includegraphics[width=#2\linewidth]{#1}}
    \notatabla{#4}
  \end{figure}}

% Nota de tabla en formato APA.
\newcommand{\notatabla}[1]{\par\vspace{2pt}{\small\raggedright\linespread{1}\selectfont\textit{Nota.} #1\par}}

\newcommand{\ruta}[1]{{\urlstyle{tt}\nolinkurl{#1}}}

% ---------- Portada (modo estudiante) ----------
\title{Construcción de una tienda en línea segura con \textit{vibe coding}: tutorial paso a paso del caso Byte}
\authorsnames{\autor}
\authorsaffiliations{{\programa}}
\course{\curso}
\professor{\docente}
\duedate{\fecha}

\abstract{Este documento es un tutorial, basado en un caso real, sobre cómo construir una aplicación web de
producción mediante \textit{vibe coding}, es decir, programando a través de instrucciones en lenguaje
natural a un agente de inteligencia artificial. El caso es Byte, una tienda en línea y punto de venta de
videojuegos, consolas y accesorios para Colombia, construida con Django 5.2, PostgreSQL 17, Celery y un
módulo de inteligencia de negocio, y desplegada en Render. Toda la construcción partió de un único
prompt maestro que fijó el rol de la inteligencia artificial, las reglas de trabajo, la arquitectura,
los requisitos funcionales y de seguridad, el plan de diez fases y los criterios de aceptación. El
tutorial explica cómo está construido ese prompt y recorre las fases mostrando, para cada una, el
fragmento del prompt que la gobernó, la forma en que se implementó y las evidencias que se generaron:
bitácoras, capturas de pantalla y resultados de las verificaciones automáticas. La aplicación terminó
con 727 pruebas, una cobertura del 93~\%, cero hallazgos de semgrep y cero fallos en OWASP ZAP baseline.
El documento deja preparada la continuación, dedicada a las pruebas de penetración sobre la aplicación
construida.}
\keywords{vibe coding, prompt maestro, desarrollo asistido por inteligencia artificial, Django, seguridad de aplicaciones, OWASP Top 10}

\begin{document}
\maketitle
\RaggedRight

% =====================================================================
% Introducción (APA: sin título "Introducción"; se repite el título del trabajo)
% =====================================================================
Construir una tienda en línea completa (catálogo, carrito, pagos, panel administrativo, inventario
multisede, punto de venta y analítica) suele requerir un equipo y varios meses. Este tutorial muestra
cómo se construyó una aplicación de ese tamaño con \textit{vibe coding}: el autor no escribió el código
línea a línea, sino que lo especificó en un prompt maestro y supervisó a un agente de inteligencia
artificial (IA) que trabaja directamente sobre el repositorio, ejecuta comandos y verifica sus propios
resultados.

El trabajo se hizo para el curso de Hacking ético. Por eso la seguridad no fue un añadido final: fue un
requisito del prompt maestro desde la primera línea, alineado con el OWASP Top 10:2025
\parencite{owasp2025}. La aplicación resultante, llamada Byte, será además el objetivo de la segunda
parte del proyecto, dedicada a las pruebas de penetración.

El documento tiene dos propósitos. El primero es enseñar un método replicable: cómo redactar un prompt
maestro, qué preparar antes de enviarlo y cómo verificar lo que la IA entrega en cada fase. El segundo
es dejar constancia de lo que ocurrió en este caso, con los fragmentos del prompt que guiaron cada fase
y las evidencias que se recogieron, incluidos los errores y su corrección.

El tutorial se organiza así. Primero se presentan el marco de referencia y el caso. Luego se describen
la preparación del entorno, la anatomía del prompt maestro, la aprobación del plan y las reglas
persistentes del proyecto. Después se recorre la construcción fase por fase y se resumen los
resultados, las decisiones que tomó la IA donde el prompt dejaba opciones y las lecciones aprendidas.
Al final se deja planteada la continuación con las pruebas de penetración. El Apéndice~A contiene el
prompt maestro completo y el Apéndice~B, un índice de las evidencias.

% =====================================================================
\section{Marco de referencia}
% =====================================================================

\subsection{Qué es el \textit{vibe coding}}

El término fue propuesto por \textcite{karpathy2025} para describir una forma de programar en la que la
persona se entrega a la conversación con un modelo de lenguaje, acepta sus cambios y casi olvida que el
código existe. En su forma más pura, el \textit{vibe coding} sirve para prototipos rápidos: se pide algo,
se prueba si funciona y se pide el siguiente cambio, sin leer en detalle lo que la IA escribió.

Para una aplicación que maneja dinero, datos personales e inventario, esa versión pura no es
suficiente. Una IA puede producir código que funciona en la demostración y aun así tiene fallas de
control de acceso, secretos expuestos o dependencias vulnerables. Por eso, en este proyecto se usó una
variante guiada por especificación: la persona no escribe el código, pero define con precisión qué debe
construirse, cómo debe trabajar la IA y cómo se demuestra que el resultado funciona.

\subsection{El prompt maestro como especificación}

El eje del método es un prompt maestro: un único mensaje extenso que cumple el papel de un documento de
requisitos y, a la vez, de un contrato de trabajo con la IA. A diferencia de una cadena de pedidos
sueltos, el prompt maestro fija desde el inicio cuatro cosas:

\begin{enumerate}[leftmargin=0.5in]
  \item \textbf{El rol y las reglas de la IA}: con qué experiencia debe actuar y qué debe hacer siempre
        (probar, analizar, documentar, hacer commits, detenerse).
  \item \textbf{El alcance}: tecnologías, arquitectura, roles, módulos y requisitos de seguridad, con un
        nivel de detalle verificable.
  \item \textbf{El plan}: diez fases con su entregable y la obligación de detenerse al final de cada una
        para la revisión humana.
  \item \textbf{La definición de terminado}: criterios de aceptación que se pueden comprobar uno a uno.
\end{enumerate}

El agente usado fue Claude Code \parencite{anthropic-claudecode}, en su aplicación de escritorio para
Windows. Este tipo de agente no solo genera texto: lee y edita archivos del proyecto, ejecuta comandos
en la terminal, levanta contenedores, corre las pruebas y opera un navegador para tomar capturas. Esa
capacidad de ejecutar y verificar es la que permite que un prompt maestro se convierta en una
aplicación funcionando.

\subsection{El caso: Byte}

Byte es la tienda en línea y punto de venta de una cadena ficticia de videojuegos, consolas y
accesorios en Colombia. Maneja precios en pesos colombianos con IVA del 19~\% incluido, envíos por
departamento y ciudad, y obligaciones legales locales: el Estatuto del Consumidor \parencite{ley1480},
que prohíbe anunciar descuentos sobre precios que no se cobraron y fija el derecho de retracto, y la ley
de protección de datos personales \parencite{ley1581}. La Tabla~\ref{tab:modulos} resume sus módulos.

\begin{table}[H]
\caption{Módulos de la aplicación Byte}\label{tab:modulos}
\small\linespread{1}\selectfont
\begin{tabularx}{\linewidth}{@{}>{\raggedright}p{3.6cm}>{\raggedright\arraybackslash}X@{}}
\toprule
Módulo & Funciones principales \\
\midrule
Tienda & Catálogo con variantes configurables y precio dinámico, búsqueda tolerante a errores, comparador, carrito de invitado, checkout con cupones y envío, pagos en modo de pruebas, factura en PDF. \\
Cuenta del cliente & Pedidos, devoluciones, retomas de consolas usadas, servicio técnico, puntos y saldo a favor, alertas de precio y disponibilidad, privacidad (descarga y eliminación de datos). \\
Panel (\ruta{/panel/}) & Resumen por rol, pedidos, inventario multisede con kardex, ajustes y traslados, compras, punto de venta con arqueo, clientes, precios, cupones, reportes, auditoría y eventos de seguridad. \\
Inteligencia de negocio & Histórico sintético, flujo de datos bronce, plata y oro, modelo estrella, segmentación RFM, pronóstico de ventas y tablero interactivo. \\
API (\ruta{/api/v1/}) & API versionada que deniega por defecto; esquema OpenAPI solo para el personal. \\
\bottomrule
\end{tabularx}
\notatabla{Elaboración propia a partir del prompt maestro y del código del repositorio.}
\end{table}

% =====================================================================
\section{Preparación del entorno}
% =====================================================================

Antes de enviar el prompt maestro conviene tener listo el entorno, porque la IA ejecuta comandos reales
y se detiene cuando falta una herramienta. La Tabla~\ref{tab:entorno} lista lo que se usó.

\begin{table}[H]
\caption{Herramientas del entorno de construcción}\label{tab:entorno}
\small\linespread{1}\selectfont
\begin{tabularx}{\linewidth}{@{}>{\raggedright}p{4.4cm}>{\raggedright\arraybackslash}X@{}}
\toprule
Herramienta & Uso en el proyecto \\
\midrule
Claude Code (aplicación de escritorio) & Agente de IA que edita el repositorio, ejecuta comandos, opera el navegador y toma capturas. \\
Windows 11 + Python 3.12 & Sistema operativo y lenguaje base. \\
uv \parencite{astral-uv} & Gestor de dependencias con archivo de bloqueo, \textit{hashes} y cuarentena de 7 días para paquetes recién publicados. \\
Docker Desktop + WSL2 & Contenedores para PostgreSQL, Valkey, el worker de Celery y la imagen de producción. \\
Git y GitHub & Control de versiones con una rama por fase y repositorio remoto para el despliegue. \\
Playwright \parencite{playwright} & Capturas de pantalla automáticas y prueba de extremo a extremo de la compra. \\
semgrep, bandit, pip-audit, Trivy y ZAP \parencite{semgrep,zap} & Análisis estático, auditoría de dependencias, escaneo de la imagen y escaneo dinámico. \\
Render \parencite{render-blueprint} & Plataforma de despliegue definida como código. \\
\bottomrule
\end{tabularx}
\notatabla{Las versiones exactas de las librerías se verificaron en PyPI y npm el día en que se fijaron y quedaron registradas en el ADR-0005 y en el archivo \ruta{uv.lock}.}
\end{table}

\begin{leccion}
Instale y pruebe la infraestructura (Docker con WSL2, base de datos, Git) antes de enviar el prompt. En
este caso Docker no estaba instalado: la IA lo detectó al intentar correr las pruebas contra
PostgreSQL, y su instalación reinició el equipo a mitad de la Fase~1. Lo que ya estaba confirmado en Git
sobrevivió; lo que solo estaba en la conversación hubo que reconstruirlo.
\end{leccion}

% =====================================================================
\section{Paso 1: redactar el prompt maestro}
% =====================================================================

El paso más importante del método es la especificación. En este caso fue un prompt de unos 28\,000
caracteres, organizado en trece secciones numeradas del 0 al 12 (Tabla~\ref{tab:anatomia}). El texto
completo está en el Apéndice~A.

\begin{table}[H]
\caption{Estructura del prompt maestro}\label{tab:anatomia}
\small\linespread{1}\selectfont
\begin{tabularx}{\linewidth}{@{}>{\raggedright}p{4.6cm}>{\raggedright\arraybackslash}X@{}}
\toprule
Sección & Para qué sirve \\
\midrule
0. Rol, contexto y reglas & Define el papel de la IA (arquitecto y desarrollador sénior con experiencia en seguridad) y las reglas no negociables: fases, commits, secretos, idioma, pruebas. \\
1. Stack tecnológico & Fija las tecnologías por capa para que la IA no improvise. \\
2. Arquitectura & Monolito modular con capa de servicios, lista de apps, identificadores públicos UUID y dinero en decimales. \\
3. Usuarios y roles & Diez roles, denegación por defecto, permisos por objeto, separación de funciones y matriz de permisos generada desde el código. \\
4--6. Tienda, panel y analítica & Requisitos funcionales de cada módulo, con el nivel de detalle de un documento de requisitos. \\
7. Seguridad OWASP 2025 & Controles obligatorios por cada categoría del Top 10, protección de formularios y suite de pruebas de seguridad. \\
8--10. API, calidad y despliegue & Reglas de la API, metas de calidad y cobertura, e infraestructura como código. \\
11. Plan de fases & Diez fases con su entregable; la IA debe detenerse al final de cada una. \\
12. Criterios de aceptación & Lista verificable de lo que debe cumplirse al terminar. \\
\bottomrule
\end{tabularx}
\notatabla{Elaboración propia a partir del prompt maestro (Apéndice~A).}
\end{table}

\subsection{Sección 0: el rol y las reglas}

La sección~0 es la que gobierna el comportamiento de la IA durante toda la obra
(Fragmento~\ref{f:reglas}). Cada regla responde a un riesgo concreto del desarrollo con IA: perder el
control del avance (fases y commits), olvidar las convenciones entre sesiones (\ruta{CLAUDE.md}),
filtrar secretos (variables de entorno), decidir de forma opaca (ADR), usar versiones inventadas u
obsoletas (verificación de versiones) y romper la lógica de negocio sin notarlo (pruebas obligatorias).

\begin{fragmento}[label=f:reglas]{Rol, contexto y reglas de trabajo}
Actúa como arquitecto de software senior y desarrollador full-stack Django, con experiencia en
e-commerce, ingeniería de datos y seguridad de aplicaciones (AppSec). Vas a construir desde cero una
aplicación web de producción para Byte, una tienda de videojuegos, consolas y accesorios en Colombia
(moneda COP, IVA 19~\%, envíos nacionales por departamento/ciudad).

Reglas de trabajo obligatorias:

Trabaja por fases (sección 11). Al final de cada fase: ejecuta los tests, el linter y el análisis de
seguridad, actualiza CHANGELOG.md, haz un commit con mensaje convencional (feat:, fix:, sec:, docs:…) y
detente para que yo revise.

Antes de escribir código, crea CLAUDE.md con las convenciones del proyecto, los comandos (correr,
testear, migrar, sembrar datos, desplegar) y las decisiones de arquitectura. Mantenlo actualizado.

Nunca pongas secretos en el código ni en el repo. Todo va por variables de entorno (.env.example
documentado, .env en .gitignore).

Si una decisión tiene varias opciones razonables, elige la más simple y segura, y documéntala como ADR
corto en docs/adr/.

Si una librería o servicio no está disponible o cambió, verifica la versión actual estable antes de
fijarla y anota el cambio.

Código, nombres de modelos y comentarios en inglés; interfaz de usuario, mensajes y documentación en
español.

Todo cambio de lógica de negocio importante debe tener tests.
\end{fragmento}

\subsection{Secciones 1 a 10: el alcance, escrito para ser verificable}

Las secciones de alcance comparten un rasgo: cada requisito se puede comprobar. El prompt no dice «la
aplicación debe ser segura», sino qué control se espera y cómo se demuestra. Los requisitos de control
de acceso del Fragmento~\ref{f:acceso} son un buen ejemplo: cada línea se tradujo después en código y en
pruebas automáticas.

\begin{fragmento}[label=f:acceso]{Requisitos de control de acceso (sección 3.2)}
Denegar por defecto: toda vista del panel exige autenticación + permiso explícito. Crea
mixins/decoradores propios (RoleRequiredMixin, ObjectPermissionMixin) y úsalos en TODAS las vistas y
endpoints.

Permisos a nivel de objeto: un cliente solo ve sus pedidos, direcciones, tickets y retomas; un gerente
solo ve las sedes asignadas.

Separación de funciones: quien crea un ajuste de inventario o un descuento > [X] \% no puede aprobarlo.

[…]

Tests que prueban, para cada endpoint sensible, que los roles sin permiso reciben 403/404 (incluye
pruebas de IDOR cambiando UUIDs).
\end{fragmento}

El prompt también deja opciones abiertas a propósito: el umbral \texttt{[X]}~\% del descuento, la
plataforma de despliegue \ruta{[PLATAFORMA_DE_DESPLIEGUE]} y varias parejas de tecnologías
(«django-guardian o servicio propio», «Wompi o Mercado Pago»). Esas opciones se resuelven con la regla
de la sección~0: elegir la más simple y segura y documentarla como ADR. La Tabla~\ref{tab:opciones}, más
adelante, muestra cómo se resolvió cada una.

\subsection{Secciones 11 y 12: el plan y la definición de terminado}

El plan de fases (Fragmento~\ref{f:fases}) divide la obra en entregables revisables, y los criterios de
aceptación permiten saber, al final, qué se cumplió y qué no.

\begin{fragmento}[label=f:fases]{Plan de fases (sección 11)}
\begin{tabular}{@{}p{4.6cm}p{9.6cm}@{}}
1. Fundaciones & Estructura, settings por entorno, Docker Compose, CI básico, CLAUDE.md, ADRs iniciales, modelo de datos y diagrama ER \\
2. Usuarios y seguridad base & Auth con allauth + MFA, axes, roles y permisos, mixins, auditoría, cabeceras/CSP, tests de acceso \\
3. Catálogo y precios & Productos, variantes y configuraciones con precio dinámico, búsqueda, filtros, comparador, historial de precios \\
4. Carrito, checkout y pagos & Carrito, cupones, envíos, pasarela sandbox, webhooks firmados, pedidos, facturas PDF, correos \\
5. Panel administrativo & Dashboard, inventario multi-sede y kardex, compras, pedidos, POS, CRM, moderación, reportes \\
6. Módulos de valor agregado & Retomas, servicio técnico, fidelidad, preventas, alertas de precio/stock, devoluciones \\
7. Big Data / BI & Generador de datos, pipeline medallón, modelo estrella, modelos analíticos, dashboards con insights, notebook PySpark \\
8. Endurecimiento de seguridad & Matriz OWASP completa, modelo de amenazas, suite de seguridad, ZAP, SBOM, revisión de hallazgos \\
9. Despliegue & Infraestructura como código, CI/CD, despliegue, datos demo, verificación end-to-end \\
10. Documentación final & README, manual de usuario por rol (con capturas), manual técnico, docs/seguridad.md, docs/big-data.md, guion de sustentación \\
\end{tabular}
\end{fragmento}

El título de la sección~11 repite la instrucción clave: «detente al final de cada una». Junto con la
regla de la sección~0 («detente para que yo revise»), convierte cada fase en una puerta de aprobación:
la IA no avanza a la siguiente hasta que la persona revisa la anterior.

\begin{leccion}
Al redactar su propio prompt maestro: (1) dé a la IA un rol con experiencia en seguridad; (2) escriba
reglas de trabajo que obliguen a probar, documentar y detenerse; (3) formule cada requisito de forma
que se pueda comprobar; (4) marque con corchetes los valores que deban decidirse y pida que cada
decisión quede en un ADR; (5) cierre con criterios de aceptación que permitan saber cuándo está
terminado.
\end{leccion}

% =====================================================================
\section{Paso 2: revisar y aprobar el plan}
% =====================================================================

Antes de escribir código conviene pedir a la IA el plan de la primera fase y revisarlo. En este caso la
IA presentó tres entregables: el árbol de carpetas, el modelo de datos en Mermaid que pedía la
sección~2 y la lista de dependencias. Cada dependencia venía con
la versión estable verificada ese día en PyPI y npm, como exigía la sección~0. Por ejemplo, propuso
Django 5.2.19, la versión LTS, en lugar de la 6.1.3, que era la más reciente, porque la sección~1 pedía
«Django 5.2 LTS (o la LTS vigente)» y la LTS tiene soporte durante todo el proyecto
\parencite{django2025}. Tras revisar el plan, la persona lo aprobó y comenzó la Fase~1.

% =====================================================================
\section{Paso 3: convertir las reglas en archivos del repositorio}
% =====================================================================

Una conversación con la IA se pierde cuando se reinicia la sesión o cuando el contexto se resume. Por
eso la sección~0 pide crear \ruta{CLAUDE.md} antes de escribir código. En este caso las reglas quedaron
en tres lugares:

\begin{enumerate}[leftmargin=0.5in]
  \item \textbf{\ruta{CLAUDE.md}}: convenciones, comandos y reglas de trabajo que el agente lee al iniciar
        cada sesión. Se actualizó en el mismo commit cada vez que cambió una convención.
  \item \textbf{Registros de decisiones de arquitectura} (ADR) en \ruta{docs/adr/}: diecisiete decisiones
        cortas con su contexto, alternativas y consecuencias.
  \item \textbf{La bitácora} en \ruta{docs/bitacora/}: un archivo por fase con las secciones \textit{Qué,
        Cómo, Por qué, Evidencias} y \textit{Pendientes}, acompañado de capturas en \ruta{docs/capturas/}.
\end{enumerate}

El fragmento siguiente, tomado de \ruta{CLAUDE.md}, muestra cómo las reglas del prompt maestro quedaron
convertidas en instrucciones permanentes.

\begin{Verbatim}
## Reglas de trabajo (obligatorias)

1. **Fases**: se trabaja por las fases de la tabla de abajo. Al cerrar cada fase
   se sigue este orden: tests, linters y análisis de seguridad → CHANGELOG.md →
   bitácora con capturas → commit convencional → **detenerse para revisión**.
2. **Ramas**: cada fase vive en `fase/NN-nombre` y se fusiona a `main` solo
   cuando el usuario la aprueba.
4. **Secretos**: nunca en el código ni en el repositorio. Todo va por variables
   de entorno (`.env.example` documentado y `.env` ignorado por git).
5. **Decisiones**: cuando haya varias opciones razonables, se elige la más
   simple y segura y se documenta como ADR en `docs/adr/`.
\end{Verbatim}

La regla «ejecuta los tests, el linter y el análisis de seguridad» quedó automatizada en
\ruta{scripts/evidencia.py}, que corre todas las verificaciones y guarda su salida como texto y como
imagen de terminal. Un segundo script, \ruta{scripts/capturas.py}, abre la aplicación con Playwright y
toma capturas en escritorio y móvil, en modo claro y oscuro, revisando además la consola del navegador
en busca de violaciones de la política de seguridad de contenido (CSP). La
Tabla~\ref{tab:verificaciones} resume las verificaciones, que son las mismas que corre el pipeline de
integración continua que pide la sección~7.

\begin{table}[H]
\caption{Verificaciones automáticas ejecutadas al cerrar cada fase}\label{tab:verificaciones}
\small\linespread{1}\selectfont
\begin{tabularx}{\linewidth}{@{}p{1.1cm}>{\raggedright}p{4.4cm}>{\raggedright\arraybackslash}X@{}}
\toprule
N.º & Verificación & Qué detecta \\
\midrule
01 & pytest con cobertura & Regresiones funcionales y de seguridad, contra PostgreSQL real. \\
02--04 & ruff, black, mypy & Errores de estilo, formato y tipos. \\
05 & bandit & Patrones inseguros en Python (SAST). \\
06 & pip-audit & Dependencias con vulnerabilidades conocidas. \\
07 & \texttt{manage.py check --deploy} & Configuración insegura de Django en producción. \\
08 & Cabeceras de \texttt{GET /} & CSP con \textit{nonce}, HSTS, X-Frame-Options, COOP, CORP y Permissions-Policy. \\
09 & Integridad de librerías JS & Que htmx, Alpine, ECharts y Tailwind coincidan con sus \textit{hashes} publicados. \\
10 & semgrep & Inyección, redirecciones abiertas, criptografía débil y otros patrones. \\
11 & Trivy & Vulnerabilidades altas o críticas en la imagen Docker. \\
12 & Matriz de permisos & Que la matriz rol × permiso documentada coincida con el código. \\
13 & Cadena de auditoría & Que la bitácora de auditoría encadenada por \textit{hash} esté íntegra. \\
14--15 & SBOM y ZAP baseline & Inventario de componentes y escaneo dinámico pasivo (desde la Fase~8). \\
\bottomrule
\end{tabularx}
\notatabla{Elaboración propia a partir de \ruta{scripts/evidencia.py} y \ruta{.github/workflows/ci.yml}.}
\end{table}

% =====================================================================
\section{Paso 4: construir por fases}
% =====================================================================

Con el plan aprobado y las reglas fijadas, la obra avanzó en las diez fases de la sección~11. Cada una
siguió el mismo ciclo (Figura~\ref{fig:ciclo}): el entregable definido en el prompt maestro, la
implementación por parte de la IA, las verificaciones automáticas, la bitácora con capturas, el commit
en la rama de la fase y la revisión humana antes de pasar a la siguiente.

\begin{figure}[H]
\caption{Ciclo de trabajo de cada fase}\label{fig:ciclo}
\centering
\begin{tikzpicture}[
  nodo/.style={draw=promptborde, fill=promptfondo, rounded corners=3pt, align=center,
               minimum width=3.2cm, minimum height=1.05cm, font=\small\linespread{1}\selectfont},
  flecha/.style={-{Stealth[length=2.2mm]}, thick, draw=promptborde}]
  \node[nodo] (p) {Entregable de la\\fase (prompt)};
  \node[nodo, right=1.1cm of p] (i) {Implementación\\por la IA};
  \node[nodo, right=1.1cm of i] (v) {Verificaciones\\automáticas};
  \node[nodo, below=1.1cm of v] (e) {Bitácora y\\capturas};
  \node[nodo, left=1.1cm of e] (c) {Commit en\\la rama de fase};
  \node[nodo, left=1.1cm of c] (r) {Revisión\\humana};
  \draw[flecha] (p) -- (i);
  \draw[flecha] (i) -- (v);
  \draw[flecha] (v) -- (e);
  \draw[flecha] (e) -- (c);
  \draw[flecha] (c) -- (r);
  \draw[flecha] (r) -- node[left, font=\footnotesize, align=right]{siguiente\\fase} (p);
  \draw[flecha, dashed] (v.north) .. controls +(0,0.9) and +(0,0.9) .. node[above, font=\footnotesize]{si algo falla, se corrige y se repite} (i.north);
\end{tikzpicture}
\notatabla{Elaboración propia a partir de la regla de trabajo de la sección~0 del prompt maestro.}
\end{figure}

La Tabla~\ref{tab:fases} resume el número de pruebas y la cobertura medidos al cerrar cada fase.

\begin{table}[H]
\caption{Pruebas y cobertura al cerrar cada fase}\label{tab:fases}
\small\linespread{1}\selectfont
\begin{tabularx}{\linewidth}{@{}p{0.6cm}>{\raggedright}X>{\RaggedLeft}p{1.5cm}>{\RaggedLeft\arraybackslash}p{1.9cm}@{}}
\toprule
N.º & Fase & Pruebas & Cobertura \\
\midrule
1 & Fundaciones & 323 & 99,8~\% \\
2 & Usuarios y seguridad base & 448 & 93,0~\% \\
3 & Catálogo y precios & 530 & 94,6~\% \\
4 & Carrito, checkout y pagos & 582 & 94~\% \\
5 & Panel administrativo & 679 & 94~\% \\
6 & Módulos de valor agregado & 699 & 93~\% \\
7 & Big data e inteligencia de negocio & 713 & 93~\% \\
8 & Endurecimiento de seguridad & 724 & 93~\% \\
9 & Despliegue & 727 & 93~\% \\
10 & Documentación final & 727 & 93~\% \\
\bottomrule
\end{tabularx}
\notatabla{Datos tomados de las bitácoras \ruta{docs/bitacora/fase-NN.md}. La Fase~9 incluye las pruebas agregadas al corregir el despliegue; la Fase~10 no cambió código. La meta del prompt maestro era una cobertura de al menos 80~\%.}
\end{table}

% ---------------------------------------------------------------------
\subsection{Fase 1: fundaciones}

\begin{fragmento}{Arquitectura (sección 2)}
Monolito modular Django con capa de servicios (la lógica de negocio NO vive en las vistas) […]

Cada app: models.py, services.py (lógica), selectors.py (consultas), forms.py, views.py, permissions.py,
tests/.

Identificadores públicos en URLs y API: UUID o slug, nunca el ID autoincremental.

Precios en DecimalField (o enteros en pesos), nunca float.

Genera el diagrama entidad-relación en Mermaid en docs/modelo-datos.md.
\end{fragmento}

\textbf{Cómo se implementó.} La IA generó veintiuna apps de Django con la misma estructura en capas y
una prueba (\ruta{tests/test_structure.py}) que impide que una app nueva rompa la convención. Configuró
los ajustes por entorno, Docker multietapa con usuario sin privilegios, Docker Compose con web, worker,
beat, PostgreSQL y Valkey, la integración continua, \ruta{CLAUDE.md}, los primeros ADR y el modelo de
datos en Mermaid. La seguridad de base (CSP con \textit{nonce}, cabeceras, cookies seguras, Argon2,
Celery con JSON, errores genéricos con identificador de correlación) se configuró desde esta fase.

La sección~1 pedía un «look “gamer” moderno» y la sección~9, un catálogo realista. La página de inicio
se construyó como la de una tienda en línea completa, con carrusel, beneficios, categorías con foto,
productos destacados, retoma, servicio técnico y marcas (Figuras~\ref{fig:inicio} y
\ref{fig:iniciomovil}). Las fotos de producto se tomaron de Wikimedia Commons, solo con licencia libre
verificable \parencite{wikimedia}: un script fija cada archivo por su \textit{hash} SHA-256, quita el
fondo de estudio, genera WebP sin metadatos y publica los créditos.

\figura{f01-inicio.jpg}{0.92}{Página de inicio de Byte en escritorio, modo oscuro}{Captura tomada con
Playwright al cerrar la Fase~1 (\ruta{docs/capturas/fase-01/inicio-escritorio-oscuro.jpg}), recortada a
la parte visible de la pantalla. Las fotos de producto provienen de Wikimedia Commons con licencia
libre.}{fig:inicio}

\figura{f01-inicio-movil.jpg}{0.36}{Página de inicio en un teléfono, modo claro}{Captura de la Fase~1
(\ruta{docs/capturas/fase-01/inicio-movil-claro.jpg}), recortada.}{fig:iniciomovil}

\textbf{Qué encontraron las verificaciones.} Al probar la imagen real en contenedores aparecieron
problemas que no se veían sin ellos. La IA los detectó con las verificaciones y los corrigió antes del
commit (Tabla~\ref{tab:hallazgos1}). La Figura~\ref{fig:resumen1} muestra el resumen de las once
verificaciones de la fase.

\begin{table}[H]
\caption{Hallazgos de la verificación con contenedores en la Fase 1}\label{tab:hallazgos1}
\small\linespread{1}\selectfont
\begin{tabularx}{\linewidth}{@{}>{\raggedright}X>{\raggedright\arraybackslash}X@{}}
\toprule
Hallazgo & Corrección \\
\midrule
Los archivos estáticos quedaban ilegibles para el usuario sin privilegios: la tienda habría salido sin estilos ni imágenes. & Permisos propios para los estáticos y prueba de regresión. \\
La imagen de producción incluía las herramientas de seguridad (978 MB). & Exportar solo las dependencias de la aplicación: la imagen bajó a 472 MB. \\
semgrep señaló el uso de SHA-1 para verificar las fotos. & Cada foto se fija por SHA-256. \\
Trivy encontró dos CVE altas en \ruta{zlib1g} de la imagen base. & La etapa base aplica las actualizaciones de seguridad de Debian. \\
El worker de Celery no registraba la ejecución de tareas. & Se aplicó la configuración de registro de Django al worker. \\
\bottomrule
\end{tabularx}
\notatabla{Resumen de la bitácora \ruta{docs/bitacora/fase-01.md}.}
\end{table}

\figura{f01-resumen.png}{0.85}{Resumen de las verificaciones automáticas de la Fase 1}{Imagen de terminal
generada por \ruta{scripts/evidencia.py} (\ruta{docs/capturas/fase-01/00-resumen.svg}). Resultado: 323
pruebas aprobadas con cobertura del 99,83~\% y cero hallazgos en los análisis de seguridad.}{fig:resumen1}

% ---------------------------------------------------------------------
\subsection{Fase 2: usuarios y seguridad base}

\begin{fragmento}{Fallas de autenticación (sección 7, A07:2025)}
django-allauth con verificación de correo, MFA TOTP (obligatorio para todo el personal), códigos de
recuperación.

django-axes: bloqueo progresivo por usuario+IP; CAPTCHA (Cloudflare Turnstile o hCaptcha) tras intentos
fallidos y en registro.

[…]

Mensajes de error genéricos (“credenciales inválidas”) para no enumerar usuarios.
\end{fragmento}

\textbf{Cómo se implementó.} La IA usó allauth para registro, verificación, MFA y sesiones, y axes para
el bloqueo progresivo por usuario e IP. Definió los roles en un solo archivo y generó desde el código la
matriz de permisos, como pedía la sección~3.2. La auditoría quedó encadenada por \textit{hash} y
protegida con un disparador de PostgreSQL que impide modificar o borrar registros. El área «Mi cuenta»
incluye el consentimiento, la descarga y la eliminación de datos que exige la ley de datos personales
\parencite{ley1581}. La Figura~\ref{fig:login} muestra el inicio de sesión y la Figura~\ref{fig:403}, la
respuesta cuando un vendedor intenta entrar a la auditoría.

\figura{f02-login.jpg}{0.85}{Inicio de sesión}{Captura de la Fase~2
(\ruta{docs/capturas/fase-02/login-escritorio-oscuro.jpg}).}{fig:login}

\figura{f02-403.jpg}{0.85}{Un vendedor intenta abrir la auditoría y recibe «Acceso denegado»}{Captura
de la Fase~2 tomada con la sesión real del usuario demo del rol vendedor
(\ruta{docs/capturas/fase-02/403-vendedor-en-auditoria-escritorio-oscuro.jpg}).}{fig:403}

\textbf{Un error de la IA, y por qué importa.} El commit de evidencias de esta fase afirmó «semgrep sin
hallazgos», pero esa ejecución había reportado tres: una posible redirección abierta en el
\textit{middleware} del personal y dos usos de SHA-1. Se detectó al revisar la salida guardada, se
corrigió en el commit siguiente y la bitácora lo registra con una nota de transparencia. La
Figura~\ref{fig:semgrep} muestra la ejecución posterior, ya sin hallazgos.

\figura{f02-semgrep.png}{0.85}{Ejecución de semgrep tras corregir los hallazgos de la Fase 2}{Imagen de
terminal de \ruta{docs/capturas/fase-02/10-semgrep.svg}: cero hallazgos en 152 reglas.}{fig:semgrep}

\begin{leccion}
La IA puede resumir mal un resultado aunque la herramienta lo haya reportado bien. Guarde siempre la
salida completa de cada verificación (no solo el resumen que escribe la IA) y contrástela antes de dar
por cerrada una fase.
\end{leccion}

% ---------------------------------------------------------------------
\subsection{Fase 3: catálogo y precios}

\begin{fragmento}{Catálogo y descubrimiento (secciones 4.1 y 4.2)}
Precio según configuración: precio base + modificadores por atributo, o precio fijo por variante; el
selector de la ficha de producto actualiza precio, disponibilidad, SKU e imagen en vivo (HTMX), y
deshabilita combinaciones que no existen.

[…]

Búsqueda full-text con tolerancia a errores (pg\_trgm) y autocompletado.
\end{fragmento}

\textbf{Cómo se implementó.} La búsqueda usa las capacidades nativas de PostgreSQL (texto completo y
similitud por trigramas) en lugar de un motor externo, para no agregar un servicio más que desplegar y
asegurar \parencite{postgresql2024}. Un ajuste fino salió de las propias pruebas: con un umbral de
similitud de 0,45, la búsqueda «timon» encontraba «verde limón»; con 0,55 dejó de hacerlo y siguió
encontrando «plaistation» (Figura~\ref{fig:busqueda}). La ficha de producto actualiza precio,
disponibilidad y SKU al cambiar la configuración y deshabilita las combinaciones que no existen
(Figura~\ref{fig:ficha}). El historial de precios es de solo inserción y el «precio más bajo en 30 días»
se calcula a partir de él, de modo que no se puede inflar un descuento, en línea con el Estatuto del
Consumidor \parencite{ley1480}.

\figura{f03-busqueda.jpg}{0.85}{Búsqueda con error de digitación: «plaistation»}{Captura de la Fase~3
(\ruta{docs/capturas/fase-03/busqueda-con-error-de-tipeo-escritorio-oscuro.jpg}).}{fig:busqueda}

\figura{f03-ficha.jpg}{0.85}{Ficha de producto con configurador de variantes}{Captura de la Fase~3
(\ruta{docs/capturas/fase-03/ficha-ps5-slim-escritorio-oscuro.jpg}), recortada.}{fig:ficha}

% ---------------------------------------------------------------------
\subsection{Fase 4: carrito, checkout y pagos}

\begin{fragmento}{Carrito y checkout (sección 4.3)}
El servidor recalcula siempre precios, descuentos, impuestos y envío; jamás confía en valores del
cliente.

[…]

Pagos con pasarela en sandbox; webhooks con verificación de firma e idempotencia; nunca almacenar datos
de tarjeta (solo tokens/referencias de la pasarela).
\end{fragmento}

\textbf{Cómo se implementó.} El pedido se crea en un único servicio transaccional que recalcula en el
servidor precios, cupón, envío e IVA, y reserva el inventario «todo o nada» bloqueando las filas en orden
para evitar interbloqueos, lo que cubre también el requisito de concurrencia de la sección~7 (A06). La
pasarela Wompi \parencite{wompi-docs} quedó detrás de la interfaz \ruta{PaymentGateway} que pedía la
sección~1, con un simulador para desarrollo que emite eventos firmados por el mismo camino que un webhook
real. El enlace de pago lleva la firma de integridad que exige la pasarela:

\begin{Verbatim}
firma = SHA-256( referencia + monto_en_centavos + "COP" + secreto_de_integridad )
\end{Verbatim}

Un pago solo se aplica con firma válida, referencia existente y monto exacto; si no, se rechaza y se
registra un evento de seguridad. Las Figuras~\ref{fig:checkout} y \ref{fig:pasarela} muestran el
checkout y la pasarela de pruebas.

\figura{f04-checkout.jpg}{0.85}{Checkout con envío, cupón y aviso de pagos de prueba}{Captura de la
Fase~4 (\ruta{docs/capturas/fase-04/checkout-escritorio-oscuro.jpg}), recortada.}{fig:checkout}

\figura{f04-pasarela.jpg}{0.85}{Pasarela de pruebas del simulador de pagos}{Captura de la Fase~4
(\ruta{docs/capturas/fase-04/pasarela-de-pruebas-escritorio-oscuro.jpg}).}{fig:pasarela}

% ---------------------------------------------------------------------
\subsection{Fase 5: panel administrativo}

\begin{fragmento}{Panel de administración (sección 5)}
Inventario: multi-sede/bodega, stock disponible/reservado/en tránsito, kardex de movimientos (entrada,
salida, ajuste, traslado, devolución) con usuario y motivo, […] ajustes con aprobación […].

POS (punto de venta): venta rápida en mostrador con búsqueda/lector de código de barras, cliente
opcional, varios medios de pago, apertura y cierre de caja con arqueo.
\end{fragmento}

\textbf{Cómo se implementó.} La IA aplicó la separación de funciones de la sección~3.2, con el umbral
\texttt{[X]} fijado en 15~\% y configurable por entorno: quien solicita un ajuste de inventario, una orden
de compra o un cupón con descuento alto no puede aprobarlo, y la base de datos lo garantiza con
restricciones \texttt{CHECK} aunque alguien salte el servicio (Figura~\ref{fig:ajustes}). Cada vista del
panel declara sus permisos y filtra por las sedes asignadas al usuario: un objeto de otra sede responde
404. El punto de venta reutiliza el mismo flujo de pedidos y reservas de la web, de modo que todas las
salidas de inventario quedan en el mismo kardex (Figura~\ref{fig:pos}).

\figura{f05-ajustes.jpg}{0.85}{Un bodeguero solicita un ajuste de inventario pero no puede aprobarlo}{Captura
de la Fase~5 (\ruta{docs/capturas/fase-05/panel-ajustes-bodega-no-aprueba-escritorio-oscuro.jpg}).}{fig:ajustes}

\figura{f05-pos.jpg}{0.85}{Punto de venta con turno de caja abierto}{Captura de la Fase~5
(\ruta{docs/capturas/fase-05/panel-pos-escritorio-claro.jpg}).}{fig:pos}

% ---------------------------------------------------------------------
\subsection{Fase 6: módulos de valor agregado}

\begin{fragmento}{Cuenta de cliente (sección 4.4)}
Retoma (trade-in): el cliente cotiza su consola/juego usado (modelo, estado, accesorios incluidos) →
obtiene oferta estimada → agenda revisión → el técnico confirma → crédito en tienda o pago.

Servicio técnico: abrir ticket de reparación o garantía, adjuntar fotos (validadas), seguir estado.
\end{fragmento}

\textbf{Cómo se implementó.} La retoma calcula una oferta estimada a partir del precio de la consola
usada y del estado declarado, y la oferta final la confirma el técnico en tienda
(Figura~\ref{fig:retomas}). El servicio técnico tiene órdenes con código de seguimiento, cotización que
aprueba el cliente y garantía (Figura~\ref{fig:servicio}). Los puntos de fidelidad y el saldo a favor se
modelaron como libros contables de solo inserción, en los que cada movimiento tiene origen, autor y
fecha, y un bloqueo por usuario evita que dos operaciones simultáneas calculen mal el saldo. Las
devoluciones respetan los plazos de retracto de la Ley~1480 y el reembolso necesita la aprobación de una
persona distinta de quien recibió el producto, según los flujos de aprobación de la sección~7 (A06).

\figura{f06-retomas.jpg}{0.85}{Cotización de una retoma desde la cuenta del cliente}{Captura de la Fase~6
(\ruta{docs/capturas/fase-06/cuenta-retomas-escritorio-oscuro.jpg}), recortada.}{fig:retomas}

\figura{f06-servicio.jpg}{0.85}{Bandeja de servicio técnico en el panel}{Captura de la Fase~6
(\ruta{docs/capturas/fase-06/panel-servicio-tecnico-escritorio-oscuro.jpg}).}{fig:servicio}

% ---------------------------------------------------------------------
\subsection{Fase 7: big data e inteligencia de negocio}

\begin{fragmento}{Pipeline de datos (sección 6.2)}
Captura de eventos: modelo Event […] registrado de forma asíncrona y sin PII cruda (ID de usuario
seudonimizado).

Bronze: extracción incremental desde PostgreSQL a archivos Parquet particionados por fecha.

Silver: limpieza, deduplicación y tipado con Polars.

Gold: modelo estrella (fact\_sales, fact\_inventory\_snapshot, fact\_events, dim\_product, dim\_customer,
dim\_date, dim\_store, dim\_channel) consultado con DuckDB.
\end{fragmento}

\textbf{Cómo se implementó.} La IA construyó el lago de datos con capas bronce, plata y oro sobre
archivos Parquet, sin servidores nuevos, y seudonimizó a los clientes antes de extraerlos, de modo que
el lago no contiene datos personales. Usó DuckDB también para la capa de plata, en lugar de Polars, y lo
documentó en el ADR-0017: con el volumen del proyecto, una sola herramienta embebida bastaba. bandit
señaló que el SQL de DuckDB se armaba con cadenas formateadas; en lugar de silenciar la alerta, la IA
reescribió las consultas con SQL constante sobre vistas registradas, como pedía la sección~7 (A05). El
tablero (Figura~\ref{fig:bi}) muestra ventas mensuales con pronóstico, categorías, sedes, canales y
segmentos RFM.

\figura{f07-bi.jpg}{0.85}{Tablero de inteligencia de negocio}{Captura de la Fase~7
(\ruta{docs/capturas/fase-07/panel-inteligencia-de-negocio-escritorio-oscuro.jpg}), recortada. Los datos
son sintéticos.}{fig:bi}

% ---------------------------------------------------------------------
\subsection{Fase 8: endurecimiento de seguridad}

\begin{fragmento}{Seguridad (sección 7)}
Genera docs/seguridad.md con una matriz de trazabilidad: categoría OWASP → riesgo en esta app → control
implementado → archivo/código → test que lo verifica. […]

En CI: bandit, semgrep --config p/django, pip-audit, manage.py check --deploy, y OWASP ZAP baseline contra
el entorno de staging o un contenedor levantado en el pipeline. El pipeline falla con hallazgos altos.
\end{fragmento}

\textbf{Cómo se implementó.} La IA escribió la matriz control por control en
\ruta{docs/seguridad/owasp-top10-2025.md}, el modelo de amenazas STRIDE con su riesgo residual y una
prueba automática por categoría OWASP. Agregó al script de evidencias la generación del SBOM CycloneDX
y un escaneo ZAP baseline \parencite{zap} contra la aplicación levantada en local. El resultado fue cero
fallos en ZAP (55 reglas superadas y 13 advertencias, todas analizadas en la bitácora), cero hallazgos de
semgrep en 414 archivos y cero vulnerabilidades altas o críticas corregibles en Trivy
(Figura~\ref{fig:resumen8}).

\figura{f08-resumen.png}{0.85}{Resumen de las quince verificaciones de la Fase 8}{Imagen de terminal
generada por \ruta{scripts/evidencia.py} (\ruta{docs/capturas/fase-08/00-resumen.svg}).}{fig:resumen8}

Conviene precisar el alcance de esta fase: son verificaciones defensivas y automáticas, hechas desde el
código y con escaneo pasivo. No sustituyen una prueba de penetración, que intenta explotar la aplicación
desde la perspectiva de un atacante. Esa es la continuación de este trabajo.

% ---------------------------------------------------------------------
\subsection{Fase 9: despliegue en Render}

\begin{fragmento}{Despliegue y disponibilidad (sección 10)}
Archivo de infraestructura como código (render.yaml Blueprint o equivalente) con: servicio web
(Gunicorn), worker de Celery, Celery Beat (o cron de la plataforma), PostgreSQL gestionado y
Redis/Key-Value gestionado.

[…]

Al final, despliega la aplicación, ejecuta seed\_demo (y un seed\_bigdata reducido acorde al plan) en
producción, verifica manualmente el flujo completo (registro → compra → panel → analítica) y entrégame
la URL pública y el usuario demo de cada rol.
\end{fragmento}

\textbf{Cómo se implementó.} La IA definió los cinco servicios en un Blueprint
\parencite{render-blueprint}: base de datos, Valkey, servicio web, worker y programador de tareas, con la
base de datos y Valkey sin acceso público. Para cumplir la regla de la sección~0 sobre secretos, todos
los secretos los genera la propia plataforma (\ruta{generateValue}), de modo que nadie los escribió ni
quedaron en el repositorio:

\begin{Verbatim}
      - key: DJANGO_SECRET_KEY
        generateValue: true
      - key: BYTE_FIELD_ENCRYPTION_KEYS
        generateValue: true # se deriva una clave Fernet (apps/core/crypto.py)
      - key: BYTE_PAYMENT_GATEWAY
        value: fake # cambiar a wompi al cargar las llaves del sandbox
      - key: BYTE_DEMO_PASSWORD
        generateValue: true # visible en el panel de Render (usuarios demo)
\end{Verbatim}

El primer despliegue falló, y los dos siguientes también. La IA leyó los registros de Render, encontró
la causa de cada fallo y lo corrigió con un commit (Tabla~\ref{tab:render}).

\begin{table}[H]
\caption{Fallos del despliegue en Render y su corrección}\label{tab:render}
\small\linespread{1}\selectfont
\begin{tabularx}{\linewidth}{@{}>{\raggedright}p{3.2cm}>{\raggedright}X>{\raggedright\arraybackslash}X@{}}
\toprule
Síntoma & Causa & Corrección \\
\midrule
El pre-despliegue salió con estado 2. & Render ejecuta el comando sin \textit{shell}: el \texttt{\&\&} llegó a \texttt{manage.py migrate} como argumento. & Envolver el comando en \texttt{sh -c}. \\
Salió con estado 127. & Render tampoco interpreta comillas: \ruta{sh} recibió la cadena con las comillas literales. & Mover los comandos a scripts en la carpeta \ruta{deploy/}. \\
Error de gunicorn al arrancar (no fatal). & Su \textit{socket} de control intentaba escribir en un directorio de solo lectura. & Desactivar el \textit{socket} de control, que no se usa. \\
\bottomrule
\end{tabularx}
\notatabla{Commits \ruta{645ae51}, \ruta{a7f82a2} y \ruta{31e5cc6} del repositorio.}
\end{table}

El script final de pre-despliegue quedó así:

\begin{Verbatim}
#!/bin/sh
# Render pre-deploy: migrations, roles and demo data (idempotent).
set -eu
python manage.py migrate --noinput
python manage.py bootstrap_demo --sin-bi
\end{Verbatim}

Antes del despliegue, la prueba de extremo a extremo con Playwright recorrió la compra completa en siete
pasos sin errores: inicio, búsqueda, ficha, carrito de invitado, inicio de sesión con fusión del carrito,
checkout, pasarela de pruebas, pago aprobado, pedido y factura (Figuras~\ref{fig:e2epago} y
\ref{fig:e2epedido}). Ya en Render, la aplicación quedó disponible en
\url{https://byte-web-v9h4.onrender.com} (Figura~\ref{fig:render}), con las cabeceras de seguridad
activas y los servicios de base de datos, Valkey, worker y programador en funcionamiento.

\figura{f09-e2e-pago.jpg}{0.8}{Pago aprobado durante la prueba de extremo a extremo}{Captura tomada por
\ruta{scripts/e2e.py} (\ruta{docs/capturas/fase-09/e2e-06-pago-aprobado.jpg}).}{fig:e2epago}

\figura{f09-e2e-pedido.jpg}{0.8}{Detalle del pedido pagado en la cuenta del cliente}{Captura tomada por
\ruta{scripts/e2e.py} (\ruta{docs/capturas/fase-09/e2e-07-pedido.jpg}).}{fig:e2epedido}

\figura{f09-render-en-vivo.jpg}{0.8}{La aplicación desplegada en Render}{Captura de
\url{https://byte-web-v9h4.onrender.com} tomada con Playwright después del despliegue.}{fig:render}

% ---------------------------------------------------------------------
\subsection{Fase 10: documentación final}

\begin{fragmento}{Entregable de la Fase 10 (sección 11)}
README, manual de usuario por rol (con capturas), manual técnico, docs/seguridad.md, docs/big-data.md,
guion de sustentación
\end{fragmento}

La última fase produjo el README, los manuales del cliente y del personal por rol, el manual técnico y
el guion de sustentación, todos en \ruta{docs/}. Como la regla de la sección~0 obligaba a documentar al
cerrar cada fase, esta fase consistió sobre todo en consolidar, no en reconstruir.

% =====================================================================
\section{Resultados}
% =====================================================================

\subsection{Pruebas y verificaciones}

La Figura~\ref{fig:pruebas} muestra cómo creció la suite de pruebas fase a fase. La cobertura se mantuvo
por encima del 93~\% desde la Fase~2, muy por encima de la meta del 80~\% que fijó el prompt maestro.

\begin{figure}[H]
\caption{Número de pruebas automáticas al cerrar cada fase}\label{fig:pruebas}
\centering
\begin{tikzpicture}
\begin{axis}[
  width=0.92\linewidth, height=6.2cm,
  ybar, bar width=14pt, ymin=0, ymax=820,
  xlabel={Fase}, ylabel={Pruebas que pasan},
  symbolic x coords={1,2,3,4,5,6,7,8,9},
  xtick=data, nodes near coords, nodes near coords style={font=\footnotesize},
  every axis plot/.append style={fill=promptborde!70, draw=promptborde},
  ymajorgrids, grid style={gray!25},
  tick label style={font=\small}, label style={font=\small}]
\addplot coordinates {(1,323) (2,448) (3,530) (4,582) (5,679) (6,699) (7,713) (8,724) (9,727)};
\end{axis}
\end{tikzpicture}
\notatabla{Elaboración propia con los datos de la Tabla~\ref{tab:fases}. La Fase~10 no agregó pruebas.}
\end{figure}

La Tabla~\ref{tab:final} consolida las cifras finales del repositorio y de las verificaciones.

\begin{table}[H]
\caption{Cifras finales del proyecto}\label{tab:final}
\small\linespread{1}\selectfont
\begin{tabularx}{\linewidth}{@{}>{\raggedright}X>{\raggedright\arraybackslash}p{6.2cm}@{}}
\toprule
Indicador & Valor \\
\midrule
Pruebas automáticas y cobertura & 727 pruebas; 93~\% de cobertura \\
Código Python (apps, configuración, scripts y pruebas) & 388 archivos; unas 25\,800 líneas \\
Apps de Django, plantillas y migraciones & 21 apps; 93 plantillas; 40 migraciones \\
Decisiones de arquitectura & 17 ADR \\
Historia del repositorio & 28 commits en 10 ramas de fase \\
semgrep & 0 hallazgos en 152 reglas y 414 archivos \\
OWASP ZAP baseline & 0 fallos; 55 reglas superadas; 13 advertencias analizadas \\
Trivy y pip-audit & 0 vulnerabilidades altas o críticas corregibles; 0 dependencias vulnerables \\
Prueba de extremo a extremo & 7 pasos sin errores \\
Despliegue & \url{https://byte-web-v9h4.onrender.com} \\
\bottomrule
\end{tabularx}
\notatabla{Conteos obtenidos con \texttt{git ls-files} y \texttt{git log} sobre la rama \ruta{main}; resultados de seguridad tomados de la bitácora de la Fase~8.}
\end{table}

\subsection{Cómo resolvió la IA las opciones abiertas del prompt}

Donde el prompt maestro ofrecía alternativas, la IA aplicó la regla «elige la más simple y segura, y
documéntala como ADR». La Tabla~\ref{tab:opciones} resume esas decisiones.

\begin{table}[H]
\caption{Opciones abiertas del prompt maestro y decisión tomada}\label{tab:opciones}
\small\linespread{1}\selectfont
\begin{tabularx}{\linewidth}{@{}>{\raggedright}p{4.4cm}>{\raggedright}p{3.4cm}>{\raggedright\arraybackslash}X@{}}
\toprule
Opción en el prompt & Decisión & Motivo (ADR) \\
\midrule
Django 5.2 LTS «o la LTS vigente» & Django 5.2.19 LTS & Soporte durante todo el proyecto (ADR-0002). \\
CSP nativa de Django o django-csp & django-csp & Django 5.2 no trae CSP nativa (ADR-0002). \\
django-guardian o servicio propio & Servicio propio & El alcance es por dueño o por sede; los \textit{querysets} filtrados deniegan por defecto sin tablas de permisos por fila (ADR-0003). \\
uv o pip-tools & uv & \textit{Hashes} y cuarentena de 7 días contra paquetes recién publicados (ADR-0005). \\
django-tailwind o CLI standalone & CLI standalone & Sin Node; menos superficie de ataque (ADR-0006). \\
Wompi o Mercado Pago & Wompi & Más directo para Colombia; queda detrás de una interfaz intercambiable (ADR-0007). \\
Redis & Valkey 8 & Compatible con Redis y es el servicio gestionado de Render (ADR-0009). \\
Turnstile o hCaptcha & Turnstile & No rastrea al usuario y casi nunca le pide resolver un reto (ADR-0012). \\
Polars para la capa de plata & DuckDB & Una sola herramienta embebida basta para el volumen del proyecto (ADR-0017). \\
Descuento > [X]~\% con aprobación & 15~\%, configurable & Variable \ruta{BYTE_HIGH_DISCOUNT_APPROVAL_PCT}. \\
\ruta{[PLATAFORMA_DE_DESPLIEGUE]} & Render & Plataforma por defecto del propio prompt. \\
\bottomrule
\end{tabularx}
\notatabla{Elaboración propia a partir de \ruta{docs/adr/}.}
\end{table}

\subsection{Criterios de aceptación}

El prompt maestro terminaba con nueve criterios de aceptación. La Tabla~\ref{tab:criterios} los
contrasta con lo logrado, sin maquillar los que quedaron parciales.

\begin{table}[H]
\caption{Criterios de aceptación del prompt maestro}\label{tab:criterios}
\small\linespread{1}\selectfont
\begin{tabularx}{\linewidth}{@{}>{\raggedright}X>{\raggedright}p{2.2cm}>{\raggedright\arraybackslash}X@{}}
\toprule
Criterio & Estado & Observación \\
\midrule
Desplegada y accesible por HTTPS con URL pública & Cumplido & En Render, con HSTS y TLS 1.2/1.3. \\
Un usuario demo por rol, cada uno con su alcance & Cumplido & Ocho usuarios; la contraseña la genera Render. \\
Configurar una consola, ver el precio cambiar, comprar y seguir el pedido & Cumplido con matiz & Se paga con el simulador; Wompi requiere cargar las llaves del sandbox. \\
Inventario sin sobreventa bajo concurrencia, con prueba & Cumplido & Reserva todo o nada con bloqueo de filas. \\
Panel con catálogo, inventario, ventas, compras, POS, clientes, retomas y reportes & Cumplido & Con alcance por sede y separación de funciones. \\
Analítica con al menos 500\,000 pedidos & Parcial & El histórico sintético tiene unas 70\,000 líneas de venta en dos años. \\
\texttt{check --deploy}, bandit, pip-audit y ZAP sin hallazgos altos & Cumplido & Ver Figura~\ref{fig:resumen8}. \\
Matriz OWASP 2025 con controles y pruebas & Cumplido & \ruta{docs/seguridad/owasp-top10-2025.md}. \\
Cobertura de al menos 80~\% y CI en verde & Cumplido con matiz & Cobertura del 93~\%; el pipeline está configurado y su ejecución en GitHub queda por confirmar. \\
\bottomrule
\end{tabularx}
\notatabla{Criterios tomados de la sección~12 del prompt maestro (Apéndice~A).}
\end{table}

% =====================================================================
\section{Lecciones aprendidas}
% =====================================================================

De la experiencia se desprenden ocho recomendaciones para quien quiera construir una aplicación seria
con \textit{vibe coding}:

\begin{enumerate}[leftmargin=0.5in]
  \item \textbf{Invierta en el prompt maestro.} Una especificación larga y concreta reemplazó decenas de
        conversaciones. Los requisitos verificables («tests que prueban 403/404 cambiando UUID») rinden
        más que los genéricos («que sea seguro»).
  \item \textbf{Trabaje por fases con una puerta humana.} La obligación de detenerse al final de cada fase
        permitió revisar y corregir el rumbo antes de que los errores se acumularan.
  \item \textbf{Verifique, no confíe.} La IA resumió mal un resultado de semgrep en la Fase~2; la salida
        completa guardada lo delató. Las pruebas de extremo a extremo y las capturas también encontraron
        errores que el código no mostraba.
  \item \textbf{Ponga la seguridad en el prompt, no al final.} Las cabeceras, la CSP, Argon2, la auditoría
        y la denegación por defecto se configuraron desde la Fase~1 porque el prompt las exigía.
  \item \textbf{Haga persistentes las reglas.} \ruta{CLAUDE.md}, los ADR y la bitácora sobrevivieron a
        reinicios y a resúmenes de contexto; la conversación no.
  \item \textbf{Automatice la evidencia.} Un script que corre todas las verificaciones y toma las capturas
        convierte el cierre de una fase en un solo comando y produce el material del informe.
  \item \textbf{Deje opciones abiertas, pero exija que se justifiquen.} Los corchetes y las alternativas
        del prompt dieron margen a la IA, y la regla de los ADR dejó cada decisión explicada.
  \item \textbf{Nunca entregue secretos por el chat.} Los secretos de producción los generó la plataforma
        y las llaves reales las carga la persona en el panel de Render.
\end{enumerate}

Este trabajo también tiene límites. La IA tomó la mayoría de las decisiones técnicas dentro del marco
fijado; una persona sin conocimientos de desarrollo y de seguridad habría tenido dificultades para
evaluar sus propuestas o detectar sus errores. Algunos criterios quedaron parciales, como el volumen de
datos de la analítica. Además, la aplicación no ha sido probada todavía por un atacante, que es
precisamente el objeto de la segunda parte.

% =====================================================================
\section{Continuación: pruebas de penetración}
% =====================================================================

Con la aplicación construida y desplegada, el proyecto continúa con una segunda parte dedicada a las
pruebas de penetración. Su propósito es contrastar los controles descritos en este tutorial con los
intentos de un atacante, siguiendo las fases habituales de un pentest (Tabla~\ref{tab:pentest}).

\begin{table}[H]
\caption{Plan de la segunda parte: pruebas de penetración}\label{tab:pentest}
\small\linespread{1}\selectfont
\begin{tabularx}{\linewidth}{@{}>{\raggedright}p{4cm}>{\raggedright\arraybackslash}X@{}}
\toprule
Fase del pentest & Objetivo previsto \\
\midrule
Alcance y reglas de enfrentamiento & Definir objetivos (la instancia de Render y una copia local con configuración de producción), límites y autorización. \\
Reconocimiento & Reunir información pública del objetivo: DNS, certificados, tecnologías y archivos expuestos. \\
Escaneo y enumeración & Identificar servicios, rutas y parámetros de la aplicación. \\
Análisis de vulnerabilidades & Buscar fallas conocidas con escáneres automáticos y revisión manual. \\
Explotación & Intentar vulnerar el control de acceso (IDOR, escalamiento de roles), la autenticación, los pagos y las entradas. \\
Post-explotación & Evaluar el impacto y verificar que la auditoría y las alertas registren los intentos. \\
Informe & Documentar hallazgos, severidad, evidencias y recomendaciones. \\
\bottomrule
\end{tabularx}
\notatabla{Elaboración propia. Los resultados de cada fase se documentarán en la segunda parte, con
capturas de la terminal como evidencia.}
\end{table}

% =====================================================================
\section{Conclusiones}
% =====================================================================

El caso Byte muestra que el \textit{vibe coding} puede producir una aplicación web de tamaño y exigencia
reales cuando parte de un prompt maestro bien construido y no de una conversación improvisada. La persona
no escribió el código, pero sí decidió el alcance, fijó las reglas, revisó cada entrega y exigió
evidencias; la IA escribió, probó, verificó y documentó.

El factor que más pesó en el resultado no fue la capacidad del modelo, sino la calidad de la
especificación y la disciplina de verificación que esa especificación imponía: cada fase cerró con las
mismas comprobaciones, y los errores, incluidos los de la propia IA, quedaron registrados y corregidos.
Ese registro es, a la vez, la mejor defensa ante una revisión y el punto de partida de la segunda parte,
en la que la aplicación será sometida a pruebas de penetración.

Por último, este documento se redactó siguiendo las normas APA en su séptima edición
\parencite{apa2020}, con apoyo en la guía de citas y referencias de \textcite{menendez2020}.

% =====================================================================
\printbibliography

% =====================================================================
\appendix

\section{Prompt maestro completo}
\label{ap:prompt}

A continuación se transcribe, sin cambios, el prompt maestro con el que comenzó la construcción. Las
tablas del original estaban separadas por tabulaciones.

{\linespread{1}\selectfont
\VerbatimInput[fontsize=\footnotesize]{prompt-maestro.txt}
}

\section{Índice de evidencias}
\label{ap:evidencias}

Todas las evidencias están versionadas en el repositorio
\url{https://github.com/NeySeb832/tienda_byte_hackingetico}. La Tabla~\ref{tab:evidencias} indica dónde
encontrar las de cada fase.

\begin{table}[H]
\caption{Ubicación de las evidencias por fase}\label{tab:evidencias}
\small\linespread{1}\selectfont
\begin{tabularx}{\linewidth}{@{}p{0.8cm}>{\raggedright}X>{\raggedright}p{4.3cm}>{\raggedright\arraybackslash}p{1.8cm}@{}}
\toprule
Fase & Bitácora & Capturas & Commit \\
\midrule
1 & \ruta{docs/bitacora/fase-01.md} & \ruta{docs/capturas/fase-01/} & \ruta{47dd77a} \\
2 & \ruta{docs/bitacora/fase-02.md} & \ruta{docs/capturas/fase-02/} & \ruta{ca3cd9b} \\
3 & \ruta{docs/bitacora/fase-03.md} & \ruta{docs/capturas/fase-03/} & \ruta{78ce40e} \\
4 & \ruta{docs/bitacora/fase-04.md} & \ruta{docs/capturas/fase-04/} & \ruta{4f692c8} \\
5 & \ruta{docs/bitacora/fase-05.md} & \ruta{docs/capturas/fase-05/} & \ruta{373a861} \\
6 & \ruta{docs/bitacora/fase-06.md} & \ruta{docs/capturas/fase-06/} & \ruta{78c400b} \\
7 & \ruta{docs/bitacora/fase-07.md} & \ruta{docs/capturas/fase-07/} & \ruta{6b81d3d} \\
8 & \ruta{docs/bitacora/fase-08.md} & \ruta{docs/capturas/fase-08/} & \ruta{66a1c2d} \\
9 & \ruta{docs/bitacora/fase-09.md} & \ruta{docs/capturas/fase-09/} & \ruta{015be4d} \\
10 & \ruta{README.md}, \ruta{docs/manuales/} & — & \ruta{c464a1e} \\
\bottomrule
\end{tabularx}
\notatabla{Cada carpeta de capturas contiene la salida de texto de cada verificación (\ruta{NN-*.txt}), su
imagen de terminal (\ruta{NN-*.svg}), el resumen (\ruta{00-resumen.svg}) y las capturas de pantalla.}
\end{table}

\end{document}
